% titel: Blatt 8, A6: Natürlicher Logarithmus
% skript: Kap. 4 (Differentialrechnung), Kap. 1.2 (Funktionalgleichung)

% Teil I: Werkzeuge, kein Lösungsweg.

Zeigen, dass eine differenzierbare Funktion $L\colon(0,\infty)\to\mathbb R$
mit den drei Eigenschaften $L(1)=0$, $L'(1)=1$, $L(xy)=L(x)+L(y)$ bereits
der natürliche Logarithmus ist.

\begin{werkzeugkasten}[Was die Funktionalgleichung wirklich bedeutet]
Die Gleichung
\begin{equation*}
  L(xy)=L(x)+L(y)\qquad\forall x,y>0
\end{equation*}
ist die zentrale \emph{Charakterisierung} des Logarithmus: Er verwandelt
Multiplikation in Addition. Genau deshalb kann man sie ableiten: Wenn
$L$ differenzierbar ist, liefert das Differenzieren \emph{einer} Variablen
(feste $x$, nach $y$) eine Aussage über $L'$. Links steht dann eine
Kettenregel ($\frac d{dy}L(xy)=L'(xy)\cdot x$, weil die innere Ableitung
von $xy$ gleich $x$ ist), rechts $L'(y)$:
\begin{equation*}
  L'(xy)\cdot x=L'(y).
\end{equation*}
Diese Identität für alle $x,y>0$ ist eine starke Zusatzbedingung; die
spezielle Wahl $y=1$ (erlaubt, weil sie für alle $y>0$ gilt) macht daraus
eine Differenzialgleichung für $L'$.

\textit{Skript:} Kap. 1.2 (Logarithmus), Kap. 4 (Ableitungen).
\end{werkzeugkasten}

\begin{werkzeugkasten}[Die Schlussfolge kompakt]
\begin{enumerate}
  \item $L(xy)=L(x)+L(y)$ nach $y$ ableiten (Kettenregel links):
        $L'(xy)\cdot x=L'(y)$.
  \item Mit $y=1$: $xL'(x)=L'(1)=1$, also $L'(x)=\frac1x$.
  \item Integration mit $L(1)=0$: Da $(\ln x)'=\frac1x$ und $\ln 1=0$,
        folgt $L(x)=\ln x$.
\end{enumerate}
Warum der dritte Schritt eindeutig ist: Zwei Stammfunktionen von
$\frac1x$ unterscheiden sich nur um eine Konstante; die Bedingung
$L(1)=0$ wählt unter ihnen genau $\ln x$.

\textit{Skript:} Kap. 4 (Ableitung von $\ln$).
\end{werkzeugkasten}

\begin{achtung}[Typische Fallen]
\begin{itemize}
  \item \textbf{Die Funktionalgleichung gilt für beide Variablen}, also
        darf man $y=1$ spezialisieren (sie gilt für alle $y>0$).
  \item \textbf{Kettenregel nicht vergessen:}
        $\frac d{dy}L(xy)=L'(xy)\cdot x$; der Faktor $x$ (die innere
        Ableitung) fehlt sonst.
  \item \textbf{$L'(1)=1$ ist die Normierung:} erst sie legt die
        Integrationskonstante fest und erzwingt $L(x)=\ln x$; ohne sie
        ginge auch $L(x)=c\ln x$ für jede Konstante $c$.
  \item \textbf{Nicht voraussetzen, dass $L=\ln$ schon bekannt ist:} wir
        \emph{zeigen} es aus den drei Axiomen.
\end{itemize}
\end{achtung}
